% Direct-authority diplomatic French for printed p.211.
\clearpage
\SourcePage{211}

\LayerTarget{tohoku:construction:complex-from-local-composition-data}
\noindent uniquement la loi de composition $(g,g')\mapsto g^{-1}g'$ dans $L$
(là où elle est définie), les intersections $D\cap g\cdot D$ ($g\in L$), et
bien entendu les groupes
$P(D\cap g_1\cdot D\cap\cdots\cap g_n\cdot D)$ (les $g_i$ dans $L$), leurs
applications de restrictions et la façon dont les $g_i\in L$ y opèrent.

\LayerTarget{tohoku:remark:1a-local-composition-computation}
\noindent\textsc{Remarques I.}\quad a) Dans les conditions précédentes,
supposons par exemple que $P$ soit le faisceau constant défini par un anneau
commutatif $k$, et que les
$D\cap g_1\cdot D\cap\cdots\cap g_n\cdot D$ soient acycliques pour $k$, et
connexes. Alors les valeurs $F_{g_1,\ldots,g_n}$ sont simplement des éléments
de $k$ (définis pour
$D\cap g_1\cdot D\cap\cdots\cap g_n\cdot D_n\ne\phi$), et le complexe de ces
cochaînes donne la cohomlogie $H^*(X;G,k)$. Cela donne une méthode explicite de
calcul si on suppose l'ensemble $L$ ci-dessus fini (ce qui implique que le
complexe $C(\mathcal U,k)^G$ est libre de type fini en toute dimension). Il
semble que ceci doive permettre par exemple la calcul de la cohomologie des
groupes modulaires à plusieurs variables (ici
$H^*(X;G,\mathbf Z)=H^*(G,\mathbf Z)$, puisque $X$ est un ouvert contractile de
l'espace euclidien), quand on connaît un domaine fondamental assez simple $D$,
(sans avoir à regarder ce qui se passe aux points multiples comme dans la
méthode [10], qui semble impraticable pour un ensemble de points fixes trop
compliqué).

\LayerTarget{tohoku:remark:1b-spectral-sequence-comparison}
b) Prenons $D=X$, donc $\mathcal U=(g\cdot X)_{g\in G}$, on constate que la
suite spectrale du théorème 5.5.3 n'est autre que la deuxième suite spectrale du
théorème 5.2.1, de terme initial $H^p(G,H^q(X,A))$. On notera d'ailleurs qu'il
y a en tous cas un homomorphisme fonctoriel canonique de la suite spectrale du
théorème 5.5.3. dans la deuxième suite spectrale du théorème 5.5.1, comme on
voit en complétant le lemme 5.5.1. resp. 5.5.2. dans ce sens.

\LayerTarget{tohoku:construction:equivariant-cech-refinement-maps}
Pour finir ce numéro, nous allons donner quelques indications sur le calcul
``céchiste'' des $H^n(X;G,A)$, généralisant les considérations de 3.8. Soient
$\mathcal U$ et $\mathcal V$ deux $G$-recouvrements, $\mathcal V$ plus fin que
$\mathcal U$, on va définir des homomorphismes canoniques

\LayerTarget{tohoku:equation:5.5.12}
\[
\tag{5.5.12}
H^n(\mathcal U;G,P)\longrightarrow H^n(\mathcal V;G,P)
\]

\LayerTarget{tohoku:proof:independence-of-equation-5.5.12}
définis pour tout $G$-préfaisceau abélien $P$, et fonctoriels en $P$. Tout
d'abord, posant $\mathcal U=(U_i)_{i\in I}$,
$\mathcal V=(V_j)_{j\in J}$, supposons qu'il existe une application
$\varphi:J\to I$ telle que $V_j\subset U_{\varphi(j)}$ pour tout $j\in J$, et
qui commute à $G$. On en déduit classiquement un homomorphisme
$\overline\varphi:C(\mathcal U,P)\to C(\mathcal V,P)$, et qui de plus ici
commute à $G$, d'où en vertu de (5.5.3) un homomorphisme (5.5.12). Ce dernier
est \textit{indépendant du choix de $\varphi$}. Si en effet $\varphi'$ est une
autre application ayant les mêmes propriétés, on construit classiquement une
homotopie bien déterminée $s$ de $\overline\varphi$ à
$\overline{\varphi'}$ :
$\overline\varphi-\overline{\varphi'}=\partial s+s\partial$, et pour des raisons
évidentes de ``transport de structure'' $s$ commute à $G$, d'où résulte que
$\overline\varphi$ et $\overline{\varphi'}$ définissent le même homomorphisme
(5.5.12).

\LayerTarget{tohoku:construction:fixed-point-free-refinement}
Si on ne suppose plus que l'on peut trouver une $\varphi$ comme ci-dessus, on
considère l'ensemble $I'=G\times I$ où on fait opérer $G$ par
$g'(g,i)=(g'g,i)$, et le $G$-recouvrement $\mathcal U'$ sans points fixes
$(U'_{(g,i)})_{(g,i)\in G\times I}$ défini par
$U'_{(g,i)}=U_{g\cdot i}=gU_i$. L'application
$\psi(g,i)=g\cdot i$ a les propriétés requises plus haut, et définit donc un
homomorphisme.

\LayerTarget{tohoku:equation:fixed-point-free-refinement-map}
\[
H^n(\mathcal U;G,P)\longrightarrow H^n(\mathcal U';G,P)
\]

\TohokuContinuation{212}
